\chapter{API Specification}

All endpoints use the versioned prefix \texttt{/api/v1/}.

\section{Endpoint Catalogue}

\begin{longtable}{@{}llll@{}}
\toprule
\textbf{Method} & \textbf{Path} & \textbf{Auth} & \textbf{Description} \\
\midrule
\endhead
GET  & \texttt{/symbols/\{name\}}              & --  & Symbol lookup \\
GET  & \texttt{/symbols/\{name\}/dependencies}  & --  & Forward deps \\
GET  & \texttt{/symbols/\{name\}/usages}         & --  & Reverse deps \\
GET  & \texttt{/search}                          & --  & Full-text search \\
GET  & \texttt{/files/\{mod\}/imports}           & --  & Module imports \\
GET  & \texttt{/viz/project/\{name\}/graph}      & --  & Viz export \\
GET  & \texttt{/charts/symbol-distribution}      & --  & Analytics \\
POST & \texttt{/admin/index/start}               & JWT & Trigger index \\
GET  & \texttt{/admin/index/status/\{id\}}       & JWT & Job progress \\
GET  & \texttt{/health}                          & --  & Readiness \\
\bottomrule
\end{longtable}

\section{Response Envelope}

Every successful response includes a \texttt{metadata} block:

\begin{lstlisting}[language={}]
{
  "metadata": {
    "indexed_at": "2025-01-15T10:30:00Z",
    "cache_hit": false,
    "freshness": { "max_age_s": 3600, "source": "database" }
  }
}
\end{lstlisting}

Error responses carry a consistent envelope with \texttt{error.code},
\texttt{error.message}, and \texttt{error.request\_id} (the correlation ID).

\section{Common Status Codes}

\begin{table}[h]
\centering
\begin{tabular}{@{}ll@{}}
\toprule
\textbf{Code} & \textbf{Meaning} \\
\midrule
200 & Success \\
202 & Accepted (async job) \\
400 & Bad request / validation \\
401 & Missing/invalid JWT \\
404 & Not found \\
429 & Rate limited \\
500 & Internal error \\
503 & Unavailable (DB down) \\
\bottomrule
\end{tabular}
\end{table}
